\documentclass[12pt]{article}
\input{iati_structure.tex}
\renewcommand{\bottomtitlespace}{3cm}

\usepackage[russian]{babel}

\begin{document}

\renewcommand{\headerLang}{Русский}
\renewcommand{\problemName}{AB13. Vision}

\renewcommand{\headerLeft}{IATI Day 1, Senior group}
\renewcommand{\headerRightFirst}{XVII INTERNATIONAL ADVANCED TOURNAMENT IN INFORMATICS}
\renewcommand{\headerRightSecond}{ONLINE 2026}

\renewcommand{\tl}{$5$ s}
\renewcommand{\ml}{$1024$ MB}
\problem{Task \problemName}
\addauthor{Author: Emil Indzhev}{2026}{04}{16}

Космический корабль USS Enterprise (с Сашкой в роли нового капитана) путешествует по вселенной, которая представляет собой бесконечную двумерную сетку секторов. В каждом секторе находится маяк, в секторе с координатами $(X, Y)$\footnote{$X$ --- номер строки (увеличивается вниз), а $Y$ --- номер столбца (увеличивается вправо).} находится маяк с \textit{силой обзора} $V_{X, Y}$. Находясь в этом секторе, корабль использует маяк, чтобы просканировать все сектора $(X', Y')$, для которых выполнено $X - V_{X, Y} \leq X' \leq X + V_{X, Y}$ и $Y - V_{X, Y} \leq Y' \leq Y + V_{X, Y}$. Таким образом, просканированные сектора образуют квадрат со стороной $2V_{X, Y} + 1$. Для каждого такого сектора корабль получает силу обзора маяка в нём, то есть $V_{X', Y'}$ (это связано с тем, что маяки работают и как передатчики, и как приёмники). После сканирования корабль телепортируется в один из увиденных секторов, поскольку маяки также работают как телепорты.

К сожалению, у маяков есть один критический недостаток: они могут отразить результат сканирования относительно оси $X$, или относительно оси $Y$, или относительно обеих сразу (то есть существует $4$ возможные ориентации результата сканирования). Заметьте, что корабль использует этот, возможно отражённый, скан, чтобы решить, в какой сектор телепортироваться, но при выполнении телепортации отражения компенсируются. Например, если скан отражен относительно оси $Y$ и корабль по скану телепортируется влево (в сторону уменьшения $Y$), то на самом деле он телепортируется вправо. Это означает, что корабль действительно телепортируется в тот сектор, который был выбран на скане.

Кроме того, у корабля нет памяти, поэтому его решение о том, куда телепортироваться, должно зависеть только от скана, созданного маяком. Дополнительно стратегия корабля должна быть детерминированной --- если на вход подан один и тот же скан, корабль обязан выбрать один и тот же сектор на скане.

Корабль начинает в неизвестных координатах и должен всегда (независимо от начального положения и независимо от того, какие ориентации сканов он получает) иметь возможность перемещаться в направлении увеличения $X$ и $Y$ сколь угодно далеко (можно считать, что требуется в итоге достичь сектора $(X, Y)$, такого что $X, Y \geq 10^{100}$). Однако для того, чтобы это было возможно, важны как стратегия перемещения корабля, так и расположение сил обзора по всей вселенной.

Именно здесь вступаете вы --- вам нужно спроектировать прямоугольный шаблон $P$ размера $N \times M$ из сил обзора, который будет бесконечно повторяться по всей вселенной: $V_{X, Y} = P_{X \bmod N, Y \bmod M}$. Вам также нужно разработать стратегию перемещения корабля, то есть функцию, которая получает на вход скан, созданный маяком, и выбирает, в какой из этих секторов телепортироваться дальше.

Поскольку маяки с большей силой обзора дороги, ваша цель --- построить корректное решение, стараясь при этом минимизировать среднюю силу обзора в шаблоне, то есть среднее значение элементов $P$.

Так как эта задача может показаться довольно сложной, вы решаете сначала рассмотреть тренировочную версию: ту же задачу, но в одномерном случае. В этой более простой задаче мы отбрасываем измерение $X$ и используем только $Y$. Тогда ваш шаблон --- это просто последовательность длины $M$, которая повторяется следующим образом: $V_Y = P_{Y \bmod M}$. Сканы, создаваемые маяками, тоже представляют собой последовательности длины $2V_Y + 1$ (содержащие сектора $Y'$, для которых выполнено $Y - V_Y \leq Y' \leq Y + V_Y$). Здесь существуют только $2$ возможные ориентации скана (обычная или отражённая), и корабль должен иметь возможность достичь $Y \geq 10^{100}$.

Ваши решения для одномерного и двумерного случаев будут оцениваться отдельно, и вы можете получить баллы за один из них, даже не имея корректного решения для другого.

\subsection{Детали реализации}

Вам нужно реализовать четыре функции: \verb|getVisionPattern1d|, \verb|getMove1d|, \verb|getVisionPattern2d| и \verb|getMove2d|. Первые две используются для одномерного случая, а вторые две --- для двумерного. В каждом тесте будет вызываться только одна из этих пар функций, но все $4$ функции должны быть реализованы (даже если реализации одной пары будут пустыми), чтобы решение считалось корректным (то есть компилировалось).

\begin{lstlisting}
 std::vector<int> getVisionPattern1d()
\end{lstlisting}

Эта функция будет вызвана один раз, если тест относится к одномерному случаю. Она должна вернуть последовательность $P$ длины $M$ --- одномерный шаблон сил обзора, который будет повторяться. Число $M$ и все элементы $P$ должны быть от $1$ до $60$.

\begin{lstlisting}
 int getMove1d(std::vector<int> v)
\end{lstlisting}

Эта функция будет вызываться один раз для каждого возможного скана, если тест относится к одномерному случаю. На вход она получает скан, созданный маяком, в виде последовательности длины $2V + 1$ (где $V$ --- сила обзора маяка в центральном секторе последовательности). Она должна вернуть, в какой сектор корабль должен телепортироваться, в виде индекса $T$ в последовательности, такого что $0 \leq T < 2V + 1$.

\begin{lstlisting}
 std::vector<std::vector<int>> getVisionPattern2d()
\end{lstlisting}

Эта функция будет вызвана один раз, если тест относится к двумерному случаю. Она должна вернуть прямоугольную таблицу $P$ размера $N \times M$ --- двумерный шаблон сил обзора, который будет повторяться. Числа $N$, $M$ и все элементы $P$ должны быть от $1$ до $60$.

\begin{lstlisting}
 std::pair<int, int> getMove2d(std::vector<std::vector<int>> v)
\end{lstlisting}

Эта функция будет вызываться один раз для каждого возможного скана, если тест относится к двумерному случаю. На вход она получает скан, созданный маяком, в виде квадратной таблицы размера $2V + 1$ на $2V + 1$ (где $V$ --- сила обзора маяка в центральном секторе квадрата). Она должна вернуть, в какой сектор корабль должен телепортироваться, в виде пары индексов $S$ и $T$ в таблице, таких что $0 \leq S, T < 2V + 1$.

Для заданной размерности $D$ ($1$ или $2$) теста проверяющая программа убедится, что ваша программа строит корректный шаблон --- то есть делает допустимые телепортационные выборы для всех возможных сканов и что ваше решение работает независимо от начальных координат и последовательности ориентаций сканов.

\subsection{Ограничения}

\vspace{0.1em}

\begin{itemize}
\item $1 \leq D \leq 2$
\item $1 \leq N, M, V_{X, Y}, V_Y \leq 60$
\end{itemize}

\newpage

\subsection{Подзадачи и система оценки}

\begin{table}[H]
\begin{tblr}{|Q[c,m]|Q[c,m]|Q[c,m]|Q[c,m]|}
\hline
\textbf{Subtask} & \textbf{Points} & \textbf{$D$} & \textbf{$A^*$} \\
\hline
$1$ & $24$ & $1$ & $1.5$ \\
\hline
$2$ & $76$ & $2$ & $1.125$ \\
\hline
\end{tblr}
\end{table}
\FloatBarrier

Ваш результат за данную подзадачу (если ваше решение для неё корректно) зависит от средней силы обзора в вашем шаблоне $P$ для неё. Пусть это значение равно $A$, а целевое значение для подзадачи равно $A^*$. Тогда ваш балл (доля баллов за подзадачу) будет вычисляться по формуле:

$$1 - {\left(1 - \frac{A^* - 1}{\max{\left(A, A^*\right)} - 1}\right)}^{0.8}$$

\subsection{Пример шаблона}

Предположим, что ваше решение возвращает следующий шаблон при $N=3$ и $M=2$:

\[
\left[
\begin{array}{c c}
1 & 2 \\
3 & 4 \\
5 & 6
\end{array}
\right]
\]

Теперь рассмотрим возможные сканы в секторе $(0, 1)$, который имеет силу обзора $2$. Существует $4$ возможные ориентации (некоторые из которых на этом шаблоне дают одинаковый скан):

\[
\begin{array}{cc}
\text{Ориентация $0$: без отражений} & \text{Ориентация $1$: отражение оси $Y$} \\
\left[
\begin{array}{ccccc}
4 & 3 & 4 & 3 & 4 \\
6 & 5 & 6 & 5 & 6 \\
2 & 1 & 2 & 1 & 2 \\
4 & 3 & 4 & 3 & 4 \\
6 & 5 & 6 & 5 & 6
\end{array}
\right]
&
\left[
\begin{array}{ccccc}
4 & 3 & 4 & 3 & 4 \\
6 & 5 & 6 & 5 & 6 \\
2 & 1 & 2 & 1 & 2 \\
4 & 3 & 4 & 3 & 4 \\
6 & 5 & 6 & 5 & 6
\end{array}
\right]
\\
\\
\text{Ориентация $2$: отражение оси $X$} & \text{Ориентация $3$: отражение обеих осей} \\
\left[
\begin{array}{ccccc}
6 & 5 & 6 & 5 & 6 \\
4 & 3 & 4 & 3 & 4 \\
2 & 1 & 2 & 1 & 2 \\
6 & 5 & 6 & 5 & 6 \\
4 & 3 & 4 & 3 & 4
\end{array}
\right]
&
\left[
\begin{array}{ccccc}
6 & 5 & 6 & 5 & 6 \\
4 & 3 & 4 & 3 & 4 \\
2 & 1 & 2 & 1 & 2 \\
6 & 5 & 6 & 5 & 6 \\
4 & 3 & 4 & 3 & 4
\end{array}
\right]
\end{array}
\]

Поскольку ориентации $0$ и $1$ совпадают, а корабль детерминирован, для них будет сделан только один вызов \verb|getMove2d|. Предположим, что ваше решение возвращает ход $(0, 1)$ на этом скане. При ориентации $0$ это означало бы перемещение на $1$ влево и на $2$ вверх, а при ориентации $1$ --- перемещение на $1$ вправо и на $2$ вверх.

\subsection{Пример грейдера}

Первым входным значением в примере грейдера является $D$. После этого он выполнит все вызовы ваших функций: сначала получит шаблон $P$, а затем запомнит все варианты телепортационных переходов. После этого он начнёт консольное взаимодействие, запрашивая начальные координаты корабля. Затем он будет выводить текущее состояние: координаты и исходный (без отражений) скан. После этого он запросит используемую ориентацию ($0$ --- без отражений, $1$ --- отразить ось $Y$, $2$ --- отразить ось $X$, $3$ --- отразить обе оси). Затем он выведет скан в заданной ориентации, а также то, какой ход выбрал корабль. Корабль будет перемещён, и процесс повторится. Обратите внимание, что пример грейдера не будет автоматически проверять корректность вашего решения.

\end{document}